\documentclass[a4paper,12pt]{article}
\usepackage[UTF8]{ctex}
\usepackage{graphicx}
\usepackage{geometry}
\usepackage{fancyhdr}
\usepackage{titlesec}
\usepackage{enumitem}
\usepackage{booktabs}
\usepackage{longtable}
\usepackage{array}
\usepackage{multirow}
\usepackage{amsmath}
\usepackage{hyperref}
\usepackage{caption}
\usepackage{setspace}

% 页面设置
\geometry{top=2.5cm,bottom=2.5cm,left=2.5cm,right=2.5cm,headheight=1.5cm,footskip=1.5cm}

% 1.5倍行距
\onehalfspacing

% 页眉页脚设置
\pagestyle{fancy}
\fancyhf{}
\fancyhead[C]{\small 软件项目开发与实践 --- 作业2}
\fancyfoot[C]{\thepage}
\renewcommand{\headrulewidth}{0.4pt}
\renewcommand{\footrulewidth}{0.4pt}

% 标题格式设置
\titleformat{\section}{\centering\zihao{-2}\heiti}{第\,\thesection\,章\quad}{0.5em}{}
\titleformat{\subsection}{\zihao{3}\heiti}{\thesubsection\quad}{0.5em}{}
\titleformat{\subsubsection}{\zihao{-3}\heiti}{\thesubsubsection\quad}{0.5em}{}

% 表格标题格式
\captionsetup{font=small,labelfont=bf,labelsep=quad}
\captionsetup[longtable]{font=small,labelfont=bf,labelsep=quad}

% 自定义三线表命令
\newcommand{\thead}[1]{\multirow{2}{*}{\textbf{#1}}}

\begin{document}

% ======================== 封面 ========================
\begin{titlepage}
\begin{center}
\vspace*{1.5cm}

{\zihao{1}\heiti 软件项目开发与实践}\\[0.8cm]
{\zihao{1}\heiti 作业2}\\[2cm]

{\zihao{2}\heiti 基于Pygame的2D RPG游戏设计与实现}\\[0.3cm]
{\zihao{3}\heiti ——以《西游记：祸起观音院》为例}\\[3cm]

\renewcommand{\arraystretch}{2.0}
\begin{tabular}{r@{\hspace{0.5em}：\hspace{0.5em}}l}
{\zihao{4}\songti 班\quad\quad 级} & {\zihao{4}\songti \underline{\makebox[120pt]{软件2401}}} \\
{\zihao{4}\songti 学\quad\quad 号} & {\zihao{4}\songti \underline{\makebox[120pt]{5120243335}}} \\
{\zihao{4}\songti 姓\quad\quad 名} & {\zihao{4}\songti \underline{\makebox[120pt]{袁杰强}}} \\
{\zihao{4}\songti 提交日期} & {\zihao{4}\songti \underline{\makebox[120pt]{2026年6月29日}}} \\
\end{tabular}

\vfill
{\zihao{4}\songti \today}
\end{center}
\end{titlepage}

% ======================== 目录 ========================
\newpage
\tableofcontents
\thispagestyle{empty}
\addtocontents{toc}{\protect\thispagestyle{empty}}

% ======================== 正文 ========================
\newpage
\setcounter{page}{1}

% ========== 第一章 ==========
\section{项目代码仓库与版本控制}

\subsection{公开Git仓库链接}

仓库地址：\url{https://github.com/YJQ-123588/home1}

\subsection{提交（Commit）历史}

\begin{longtable}{|c|l|l|}
\hline
\textbf{序号} & \textbf{提交信息} & \textbf{提交时间} \\
\hline
\endfirsthead
\hline
\textbf{序号} & \textbf{提交信息} & \textbf{提交时间} \\
\hline
\endhead
\hline
\endfoot
1 & 实践课-西游记-初始1.0 & 2026-06-24 03:29:30 \\
\hline
2 & Create README.md & 2026-06-24 06:22:32 \\
\hline
3 & Update README.md & 2026-06-24 06:24:15 \\
\hline
4 & 西游记2.0 & 2026-06-24 06:37:52 \\
\hline
5 & Merge branch 'main' 合并远端仓库更新 & 2026-06-24 06:44:41 \\
\hline
6 & main分支同时存在几套项目源码 & 2026-06-24 07:37:38 \\
\hline
7 & 新增PythonProject3-pygame,西游记3.0 & 2026-06-25 00:12:45 \\
\hline
8 & 仓库同步1.0/2.0/3.0三套源码 & 2026-06-25 01:12:38 \\
\hline
9 & 新增PythonProject-pygame3.0 & 2026-06-25 01:45:36 \\
\hline
10 & 新增PythonProject-pygame3.1 & 2026-06-25 02:37:57 \\
\hline
\end{longtable}

\subsection{版本控制反思}

本次项目采用增量式开发策略，按照功能模块逐步提交：

\begin{enumerate}[label=(\arabic*)]
    \item \textbf{初始1.0版本}：完成基础游戏框架，包含玩家移动、场景切换等核心功能；
    \item \textbf{2.0版本}：新增战斗系统、NPC对话、存档功能；
    \item \textbf{3.0版本}：完善战斗逻辑、增加避战选项、优化地图碰撞；
    \item \textbf{3.1版本}：修复HP上限bug、完善存档系统、增加状态栏UI。
\end{enumerate}

每次提交都对应一个完整的功能模块，便于追溯和回滚。遇到Bug时单独提交修复，确保版本稳定性。

% ========== 第二章 ==========
\section{AI协作能力专项测评}

\subsection{与AI结对编程——调试/功能实现实录}

\subsubsection{问题描述}

在开发战斗系统时，遇到了战斗状态机逻辑混乱的问题。具体表现为：

\begin{enumerate}[label=(\arabic*)]
    \item 玩家攻击后，怪物反击动画不播放；
    \item 战斗结束条件判断错误，导致无限循环；
    \item 大招解锁状态在场景切换后丢失。
\end{enumerate}

\subsubsection{AI提示词（Prompt）}

\begin{quote}
\textit{``我在开发一个Pygame RPG游戏的战斗系统，使用有限状态机管理战斗流程。当前问题是：}
\begin{enumerate}[label=(\arabic*)]
    \item \textit{STATE\_PLAYER\_ATK状态执行完后，没有正确切换到STATE\_MONSTER\_HURT}
    \item \textit{战斗结束判断条件有问题，击败怪物后无法正常退出战斗}
    \item \textit{ultimate\_unlocked变量在场景切换后被重置}
\end{enumerate}
\textit{请帮我分析fight\_scene.py中的状态机逻辑，找出问题所在并提供修复方案。''}
\end{quote}

\subsubsection{AI给出的解决方案}

\begin{enumerate}[label=(\arabic*)]
    \item 检查状态转换条件，确保每个状态都有明确的next\_state；
    \item 添加战斗结束标志，在怪物HP$\leq$0时正确设置fight\_over=True；
    \item 将ultimate\_unlocked状态同步到Game主控类，实现场景间持久化。
\end{enumerate}

\subsubsection{最终处理方式}

选择 \textbf{B}：理解思路后，手动重写了代码并进行了适配。

\begin{itemize}
    \item 重新设计了状态机的状态转换表；
    \item 添加了状态持续时间控制，避免状态跳转过快；
    \item 将关键状态变量同步到Game类实现持久化。
\end{itemize}

\textbf{学到的最重要一课}：AI能够快速定位问题并提供解决思路，但需要开发者理解底层逻辑后进行适配，而不是盲目复制代码。

% ========== 第三章 ==========
\section{架构优化}

\subsection{系统概要类图（优化后）}

\begin{verbatim}
+-------------------------------------------+
|              Game (主控类)                 |
+-------------------------------------------+
| - state: int                              |
| - player_hp: int                          |
| - player_max_hp: int                      |
| - player_atk: int                         |
| - ultimate_unlocked: bool                 |
| - temple_defeated: list                   |
| - temple_boss_appeared: bool              |
+-------------------------------------------+
| + change_state(new_state: int)            |
| + get_game_state() -> dict                |
| + save_game(slot: int)                    |
| + load_game(slot: int)                    |
+-------------------------------------------+
                    |
                    | 管理
                    v
+-------------------------------------------+
|            Scene (场景基类)                |
+-------------------------------------------+
| + update(events, dt)                      |
| + draw(surface)                           |
+-------------------------------------------+
          |
    +-----+------+----------+----------+
    v            v          v          v
+--------+  +--------+  +--------+  +--------+
| Start  |  |Village |  |Temple  |  | Fight  |
| Scene  |  | Scene  |  | Scene  |  | Scene  |
+--------+  +--------+  +--------+  +--------+
                    |
                    | 包含
                    v
+-------------------------------------------+
|           Character (角色类)              |
+-------------------------------------------+
| - hp: int                                 |
| - max_hp: int                             |
| - atk: int                                |
+-------------------------------------------+
| + take_damage(amount: int)                |
| + is_alive() -> bool                      |
+-------------------------------------------+
          |
    +-----+-----+
    v            v
+--------+  +--------+
| Player |  |  NPC   |
| (悟空) |  | (农民) |
+--------+  +--------+
\end{verbatim}

\subsection{优化分析}

\subsubsection{优化原因}

\begin{enumerate}[label=(\arabic*)]
    \item \textbf{单一职责原则}：Game类只负责状态管理，场景类只负责渲染和交互；
    \item \textbf{组合模式}：场景包含角色，角色包含动画；
    \item \textbf{场景模板方法}：基类定义流程，子类实现具体逻辑；
    \item \textbf{状态持久化}：关键数据同步到Game类，实现跨场景保存。
\end{enumerate}

\subsubsection{采用的设计模式}

\begin{enumerate}[label=(\arabic*)]
    \item \textbf{单例模式}：Game类作为全局唯一状态管理器；
    \item \textbf{观察者模式}：状态变化时通知相关组件更新；
    \item \textbf{状态模式}：战斗系统使用有限状态机管理流程；
    \item \textbf{模板方法模式}：场景基类定义渲染流程，子类重写具体实现。
\end{enumerate}

\end{document}
